\documentclass[11pt]{article}
\usepackage[T1]{fontenc}
\usepackage[utf8]{inputenc}
\usepackage{lmodern,microtype,amsmath,amssymb,amsthm,mathtools,booktabs,array,longtable}
\usepackage[margin=1in]{geometry}
\usepackage{xcolor,hyperref,fancyhdr}
\hypersetup{colorlinks=true,linkcolor=blue!45!black,urlcolor=blue!45!black,citecolor=blue!45!black,pdftitle={A Pearson determinant bound below five for the irrationality measure of zeta(3)},pdfauthor={Zeta Irrationality project}}
\newtheorem{theorem}{Theorem}[section]
\newtheorem{lemma}[theorem]{Lemma}
\newtheorem{proposition}[theorem]{Proposition}
\newtheorem{corollary}[theorem]{Corollary}
\theoremstyle{definition}\newtheorem{definition}[theorem]{Definition}
\theoremstyle{remark}\newtheorem{remark}[theorem]{Remark}
\DeclareMathOperator{\arsinh}{arsinh}
\DeclareMathOperator{\arcosh}{arcosh}
\DeclareMathOperator{\artanh}{artanh}
\DeclareMathOperator{\diag}{diag}
\DeclareMathOperator{\rank}{rank}
\DeclareMathOperator{\nullity}{nullity}
\DeclareMathOperator{\even}{even}
\newcommand{\dd}{\,\mathrm d}
\newcommand{\pos}[1]{(#1)_+}
\newcommand{\J}{\mathcal J_{3,X}}
\newcommand{\Bfun}{\mathcal B}
\newcommand{\Q}{\mathbb Q}
\newcommand{\Z}{\mathbb Z}
\newcommand{\R}{\mathbb R}
\newcommand{\Fp}{\mathbb F_p}
\pagestyle{fancy}\fancyhf{}\fancyhead[L]{\small Zeta Irrationality: Pearson determinants}\fancyhead[R]{\small $\mu(\zeta(3))<4.974$}\fancyfoot[C]{\thepage}
\setlength{\headheight}{14pt}
\setlength{\parskip}{4pt}
\setlength{\emergencystretch}{2em}
\title{A Pearson determinant bound below five\\for the irrationality measure of $\zeta(3)$}
\author{Zeta Irrationality project\\\small Research report and reproducible certificates}
\date{18 November 2026}
\begin{document}
\maketitle
\begin{abstract}
We apply the project's growing determinant and quantitative response method to the positive cubic functional underlying Pr\'evost's remainder Pad\'e construction. Two unequal factorial windows, with endpoint ratios $2/25$ and $8/25$ and matrix-size ratio $51/50$, give
\[
 \mu(\zeta(3))\le
 \frac{2225439040605380000}{447428816560392051}
 =4.9738393197680857\ldots<4.974<5.
\]
The proof combines a cubic local-divisibility calculation, an explicit weighted-polynomial decay bound, a distributed-zero compression of a positive contour matrix, and a Pearson bound of two on determinant-root multiplicities. The arithmetic profile is integrated over 6,752 exact rational affine cells. The analytic inequalities use directed interval arithmetic, rational angular charges, and monotone Riemann sums. The decisive certified response margin is $0.004682016878632102$. This is an ordinary mathematical proof with finite numerical certificates, not a formal proof-assistant verification. The construction and all numerical inputs needed to replay the finite checks are supplied with this report.
\end{abstract}
\clearpage
\tableofcontents
\newpage
\section{Result, construction, and scope}
For a real irrational $\xi$, its irrationality exponent is
\[
 \mu(\xi)=\sup\{\nu:\ |\xi-p/q|<q^{-\nu}
       \text{ for infinitely many coprime }p\in\mathbb Z,q\in\mathbb N\}.
\]
A smaller upper bound is stronger. The result of this investigation is the following.
\begin{theorem}\label{thm:main}
The cubic Pearson determinant construction described below satisfies
\[
 \boxed{\displaystyle
 \mu(\zeta(3))\le
 \frac{2225439040605380000}{447428816560392051}
 <\frac{2487}{500}=4.974.}
\]
In particular $\mu(\zeta(3))<5$.
\end{theorem}
The exponent in the theorem is an upper bound, not an assertion that the actual irrationality exponent equals it. The argument does not supply an explicit denominator threshold $q_0$, and no optimality claim for this determinant family is made.

\subsection{The precise determinant}
Put $\xi=\zeta(3)$ and let $\mathcal J_{3,X}$ be the rational functional determined by
\begin{align}
 \mathcal J_{3,X}(t^k)&=\frac{(-1)^k(2k+3)B_{2k+2}}2,\label{eq:moments}\\
 \mathcal J_{3,X}\!\left(\frac1{t+i^2}\right)
 &=i^2X-i^2H_{i-1}^{(3)}-\frac1{2i}-\frac12,
 \qquad i\ge1.\label{eq:poles}
\end{align}
Here $B_1=-1/2$ and $H_j^{(3)}=\sum_{a=1}^j a^{-3}$. Polynomial division and partial fractions extend the functional to all rational functions used below. At $X=\xi$ it is integration against the positive measure
\[
 \dd\mu_3(x^2)=x^3\frac{\dd^2}{\dd x^2}
                  \frac1{e^{2\pi x}-1}\dd x
 =4\pi^2x^3\sum_{k\ge1}k^2e^{-2\pi kx}\dd x,
 \qquad x>0.
\]
The polynomial moments follow by two integrations by parts; the simple-pole formula is the cubic Stieltjes identity. This is the Pr\'evost-type scalar input. Our quantitative construction uses determinants of growing size.

Define
\[
 C_m(t)=\prod_{a=0}^{m-1}(t+a^2),\quad
 D_N(t)=\prod_{i=1}^N(t+i^2),\quad
 b_j(t)=\frac{(-1)^j C_j(t)}{(2j)!}.
\]
Take $N$ through positive multiples of $50$, and set
\[
 m_1=\frac{2N}{25},\qquad m_2=\frac{8N}{25},\qquad r=\frac{51N}{50}.
\]
The factorial normalization is
\begin{equation}\label{eq:FN}
 F_N(X)=\det\left[(N!)^2\mathcal J_{3,X}
 \left(\frac{b_{m_1+j}(t)b_{m_2+k}(t)}{D_N(t)}\right)
 \right]_{0\le j,k<r}.
\end{equation}
An invertible constant row and column change relates this determinant to
\begin{equation}\label{eq:Hmain}
 H_N(X)=\left[\mathcal J_{3,X}
 \left(\frac{C_{m_1}(t)C_{m_2}(t)t^{j+k}}{D_N(t)}\right)
 \right]_{0\le j,k<r}.
\end{equation}
At $\xi$, the latter is a positive definite Gram matrix. Let $P_N\in\mathbb Z[X]$ be the primitive integer polynomial proportional to $F_N$. Multiplying by a rational constant changes neither its zeros nor its normalized response.

\subsection{The three quantitative inputs}
The following constants will be established and certified:
\begin{align}
 \deg P_N&=\frac{17N}{25}+1,\label{eq:degree-main}\\
 \log|P_N(\xi)|&\le-\varepsilon N^2+o(N^2),
 &\varepsilon&\ge0.961102592443632102,\label{eq:small-main}\\
 I(T):=\int_0^T g(t)\dd t&\le0.956420575565,
 &T&=\frac{17}{125}.\label{eq:response-main}
\end{align}
Every finite zero of $P_N$ has multiplicity at most two. The response envelope $g$ is positive and decreasing, and
\[
 6.545408942519\le g(T)\le6.545408942957.
\]
The gap between \eqref{eq:small-main} and \eqref{eq:response-main}, rather than the small value alone, proves the irrationality-measure bound.

\subsection{What is new relative to the scalar specialization}
The positive functional and its scalar Pad\'e approximants have historical antecedents in Pr\'evost's work. We do not claim a new scalar starting point. The useful extra freedom is the combination of two unequal windows, a growing determinant, prime-by-prime determinant divisibility, and a response estimate on a large polynomial subspace. In particular the asymmetry reduces the degree rate to $\lambda=1-s_2=17/25$ while leaving sufficient analytic decay. The proof below uses the earlier distributed-zero compression; a stronger Green-function refinement explored during this investigation is unnecessary for Theorem~\ref{thm:main}.

\section{Certified numerical inputs and the final deduction}\label{sec:certificate}
The arithmetic and analytic sections below define $C_{\rm ar}$ and $C_{\rm an}$ and prove
\[
 \varepsilon=C_{\rm an}^{\rm low}-C_{\rm ar}^{\rm up}.
\]
All displayed terminating decimals in the following table are outward bounds.
\begin{center}
\begin{tabular}{@{}p{.45\linewidth}r@{}}
\toprule
Quantity & Certified value or bound\\
\midrule
Window endpoints $(s_1,s_2)$ & $(2/25,8/25)$\\
Matrix-size ratio $\delta$ & $51/50$\\
Degree rate $\lambda$ & $17/25$\\
Contour radius parameter $u$ & $1000001/1000000$\\
Spectral cutoff $T$ & $17/125$\\
Auxiliary cutoff $T_{\rm cap}$ & $7/50$\\
Arithmetic cost $C_{\rm ar}$ & $\le3.852564338452367898$\\
Analytic decay $C_{\rm an}$ & $\ge4.813666930896$\\
Net determinant decay $\varepsilon$ & $\ge0.961102592443632102$\\
Integrated simple response $I(T)$ & $\le0.956420575565$\\
Response slack $S=\varepsilon-I(T)$ & $\ge0.004682016878632102$\\
Endpoint $a=g(T)$ & $\le6.545408942957$\\
\bottomrule
\end{tabular}
\end{center}
The exact arithmetic bound is
\[
 C_{\rm ar}^{\rm up}
 =\frac{1926282169226183949}{500000000000000000}.
\]
The analytic lower bound is $300854183181/62500000000$, and the simple response upper bound is $191284115113/200000000000$. Consequently
\[
 S\ge S_0:=\frac{2341008439316051}{500000000000000000}>0.
\]
The response transfer theorem proved below gives
\[
 \mu(\xi)\le\frac{\lambda a}{S+Ta}.
\]
For $S>0$, this expression increases with $a$ and decreases with $S$. Substitution of $a\le6545408942957/10^{12}$ and $S\ge S_0$ therefore gives exactly
\[
 \mu(\xi)\le
 \frac{2225439040605380000}{447428816560392051}
 =4.9738393197680857065\ldots<4.974.
\]
The coarser deduction $\mu(\xi)<\lambda/T=5$ requires only the strict inequality $I(T)<\varepsilon$. The stated $4.974$ is obtained from the same certificate by retaining its positive slack.

\subsection{How the finite checks work}
The arithmetic integral uses the cutoff $\beta=p/N\ge1/500$. Every floor breakpoint and every change of a minimum, maximum, or ordering is split at its exact rational location. The profile is affine on each of 6,752 accepted cells, whose union is $[1/500,67/25]$. Each cell integral is rounded upward to a multiple of $10^{-18}$. The accumulated rounding allowance is at most $6.752\cdot10^{-15}$. The remaining prime tail has the proved upper bound $4\delta/500=51/6250$.

The analytic check uses 45-decimal-digit directed intervals. Ordinary floating-point root solvers only propose rational brackets; each accepted bracket is checked by strict interval signs before it enters a proof bound. The calculation uses:
\begin{itemize}
\item 8,192 left Riemann panels for the increasing analytic-decay integrand;
\item 2,048 left Riemann panels for a positive angular-deficit coefficient;
\item 128 rational midpoint angular charges, justified by convexity;
\item 1,024 left Riemann panels for the decreasing simple response.
\end{itemize}
The selected coefficient is
\[
 c_\theta=\frac{3835733171}{31250000000}=0.122743461472.
\]
The angular endpoint at $T_{\rm cap}$ is enclosed by
\[
 0.318106205460\le B(T_{\rm cap})\le0.318106205481<1.
\]
The physical support endpoint satisfies
\[
 0.911737892141\le b(\delta)\le0.911737892162<u.
\]
These strict guards are well separated from their boundaries. The interval checker records all final bounds and the number of verified root brackets in its JSON output. A separate exact-cell review verified coverage, each area rounding, the sum, and the SHA256 hash of the arithmetic cell file.

% Standalone analytic and determinant core for the zeta(3) measure bound.
% Requires only amsmath and amssymb; no local theorem environments or macros.
% The coefficient-clearing estimate is the input from the separate arithmetic note.

\section{The positive determinant and its analytic estimates}
\label{z3a:section}

This section proves the analytic and determinant inputs for the bound
$\mu(\zeta(3))<4.974$.  Only the coefficient-clearing estimate stated below is
delegated to the arithmetic section.  In particular, the polynomial norm
comparison and the conversion from determinant values to an irrationality
measure are proved here.  No modified physical equilibrium or Green-function
improvement is needed.

\subsection{Positive measure, rational functional, and the two windows}

Put $f(x)=(e^{2\pi x}-1)^{-1}$ and define the positive measure $\mu_3$ by
\[
 \int h(t)\,\mathrm d\mu_3(t)
 =\int_0^\infty h(x^2)x^3f''(x)\,\mathrm dx.
\]
Its density in the $x$ variable is strictly positive, tends to $1/\pi$ at
zero, and is bounded above by $C(1+x^3)e^{-2\pi x}$ and below by
$c e^{-2\pi x}$, for absolute positive constants $c,C$.  These assertions
follow from the Laurent expansion at zero and the positive series
$f''(x)=(2\pi)^2\sum_{n\ge1}n^2e^{-2\pi nx}$.
With $B_1=-1/2$, direct integration gives
\[
 \int t^j\,\mathrm d\mu_3(t)
 =\frac{(-1)^j(2j+3)}2 B_{2j+2}.
\]
The Stieltjes transform is
\[
 \int\frac{\mathrm d\mu_3(t)}{t+y^2}
 =y^2\zeta(3,y)-\frac1{2y}-\frac12\qquad(y>0).
\]
For example, if $K(y)=\int_0^\infty xf(x)/(x^2+y^2)\,\mathrm dx$,
Binet's identity says $2K(y)=\log y-(2y)^{-1}-\psi(y)$.
Two integrations by parts make the displayed Stieltjes integral
$y^2K''(y)$, which is the claimed expression.  The boundary terms vanish
by the same Laurent expansion and exponential decay.

Consequently the following functional on rational functions with simple
poles among the negative integer squares is well defined by partial fractions:
\begin{align*}
 \mathcal J_X(t^j)&=\frac{(-1)^j(2j+3)}2 B_{2j+2},\\
 \mathcal J_X((t+i^2)^{-1})
 &=i^2X-i^2H_{i-1}^{(3)}-\frac1{2i}-\frac12.
\end{align*}
At $\xi=\zeta(3)$ it is integration against $\mu_3$.
Write
\[
 C_m(t)=\prod_{a=0}^{m-1}(t+a^2),\qquad
 D_N(t)=\prod_{i=1}^N(t+i^2),\qquad
 B_j(t)=\frac{(-1)^jC_j(t)}{(2j)!}.
\]
Throughout, $N$ runs through multiples of $50$, and
\[
 s_1=\frac2{25},\quad s_2=\frac8{25},\quad
 \delta=\frac{51}{50},\qquad m_i=s_iN,\quad r=\delta N.
\]
The factorial-window determinant is
\[
 F_N(X)=\det\left[(N!)^2\mathcal J_X
       \left(\frac{B_{m_1+j}B_{m_2+k}}{D_N}\right)
       \right]_{0\le j,k<r}.
\]
Its sign at $\xi$ is fixed and can be removed.  Constant invertible row and
column changes identify it with the symmetric Hankel pencil
\[
 H_N(X)=\left[\mathcal J_X
 \left(\frac{C_{m_1}C_{m_2}t^{j+k}}{D_N}\right)
 \right]_{0\le j,k<r}.
\]
The matrix $G_N=H_N(\xi)$ is positive definite.  If $M_N=[X]H_N(X)$ and
$T_N=G_N^{-1/2}M_NG_N^{-1/2}$, every nonzero rational normalization $P_N$
of this determinant satisfies
\begin{equation}\label{z3a:ratio}
 P_N(\xi+h)=P_N(\xi)\det(I+hT_N).
\end{equation}

The coefficient matrix $M_N$ is a sum of evaluation rank-one matrices at
$-i^2$, $m_2\le i\le N$, with nonzero weights.  Evaluation at these distinct
nodes is onto because $r>N-m_2+1$.  Thus
$\operatorname{rank}M_N=N-m_2+1$.  Since $T_N$ is real symmetric, its nonzero
eigenvalues count the degree of $\det(I+hT_N)$, and
\begin{equation}\label{z3a:degree}
 \deg P_N=N-m_2+1=\lambda N+1,\qquad \lambda=\frac{17}{25}.
\end{equation}

The arithmetic input is the existence of positive integers $d_N$ such that
$d_NF_N\in\mathbb Z[X]$ and
\begin{equation}\label{z3a:arithmetic-input}
 \limsup_{N\to\infty}\frac{\log d_N}{N^2}
 \le C_{\rm ar}:=
 \frac{1926282169226183949}{500000000000000000}.
\end{equation}
One may take $P_N$ primitive, changing its sign so $P_N(\xi)>0$.
Dividing $d_NF_N$ by its integer content only improves the small-value bound.

\subsection{An explicit equilibrium measure}

For $a\ge0$ set $I_a(y)=\int_0^a\log(y^2+x^2)\,\mathrm dx$ and
\[
 V(y)=2\pi|y|+I_1(y)-I_{s_1}(y)-I_{s_2}(y).
\]
For $b,s>0$, let $\nu_{b,s}$ be the measure on $[-b,b]$ with density
\[
 \frac{s\sqrt{b^2-x^2}}
 {2\pi\sqrt{s^2+b^2}(x^2+s^2)}.
\]
Its Cauchy transform and mass are, respectively,
\[
 \int\frac{\mathrm d\nu_{b,s}(x)}{z-x}
 =\frac{z-s\sqrt{z^2-b^2}/\sqrt{s^2+b^2}}{2(z^2+s^2)},
 \qquad \frac12\left(1-\frac{s}{\sqrt{s^2+b^2}}\right).
\]
The square root has the branch asymptotic to $z$ at infinity.  The formula
can be verified by its jump across the interval, cancellation of its
apparent poles, and its behavior at infinity.  Define
\[
 c(s)=\mathbf1_{s\ge s_1}+\mathbf1_{s\ge s_2}+\mathbf1_{s\le1},
 \qquad \eta_b=\int_0^\infty c(s)\nu_{b,s}\,\mathrm ds.
\]
Here $1\le c(s)\le3$.  In particular $\eta_b$ is positive.  Its mass is
\begin{equation}\label{z3a:mass}
 v(b)=\frac{b+1-s_1-s_2-\sqrt{1+b^2}
             +\sqrt{s_1^2+b^2}+\sqrt{s_2^2+b^2}}2.
\end{equation}
This is strictly increasing from zero to infinity; indeed $v'(b)\ge1/2$.
Write $b(v)$ for the inverse and $U_b$ for the logarithmic potential of
$\eta_b$.  The Cauchy transforms give
\begin{equation}\label{z3a:equilibrium}
 4U_b(y)-V(y)=\ell(v)\quad (|y|\le b),\qquad
 4U_b(y)-V(y)\le\ell(v)\quad (y\in\mathbb R).
\end{equation}
For completeness, outside the support the derivative for $y>b$ is
\[
 -2\sqrt{y^2-b^2}\int_0^\infty
 \frac{c(s)s}{\sqrt{s^2+b^2}(y^2+s^2)}\,\mathrm ds<0;
\]
on the interior it is zero.  Continuity of the logarithmic potential
extends the contact identity to the endpoints.  The contact constant is
\begin{align}
 \ell(v)={}&2+2\sum_{i=1}^2s_i\log s_i-2\sum_{i=1}^2s_i
       +4v\bigl(\log(b/2)-1\bigr)\notag\\
 &-2\log\frac{1+\sqrt{1+b^2}}2
       +2\sum_{i=1}^2s_i
           \log\frac{s_i+\sqrt{s_i^2+b^2}}{2s_i},
 \qquad b=b(v).\label{z3a:ell}
\end{align}
One way to verify the constant, without evaluating a singular integral,
is to differentiate the displayed densities: $\partial_v\eta_{b(v)}$
is the unit arcsine measure on $[-b,b]$.  Its potential on that interval
is $\log(b/2)$, so $\ell'(v)=4\log(b/2)$ and
$\ell(0)=-V(0)$.  These determine \eqref{z3a:ell}.

All potentials used here are continuous.  Indeed, if $\Lambda_b$ has
density $\pi^{-1}\operatorname{arcosh}(b/|x|)$, direct integration gives
$\Lambda_b=2\int_0^\infty\nu_{b,s}\,\mathrm ds$; hence
$\frac12\Lambda_b\le\eta_b\le\frac32\Lambda_b$.
After the map $x\mapsto x^2$, the density is bounded on compact parameter
ranges by $Cy^{-1/2}(1+|\log y|)$.  These bounds also justify all potential
derivatives by first truncating the logarithmic kernel.

\subsection{Small determinant values}

Let $\widehat h_{k,N}$ be the minimum squared norm of a monic degree-$k$
polynomial for $C_{m_1}C_{m_2}\,\mathrm d\mu_3/D_N$.  Triangular changes to
monic orthogonal polynomials on both sides give the exact identity
\begin{equation}\label{z3a:norm-product}
 |F_N(\xi)|=\prod_{k=0}^{r-1}
 \frac{(N!)^2}{(2m_1+2k)!(2m_2+2k)!}\,\widehat h_{k,N}.
\end{equation}
The different windows cause no difficulty: each side has one polynomial
of every degree $0,\ldots,r-1$, and the two leading coefficients supply
the two factorials in \eqref{z3a:norm-product}.

Here is an explicit trial-polynomial argument.  For $v=k/N$, push $\eta_{b(v)}$
forward by $x\mapsto x^2$ and partition it into $k$ consecutive cells of
mass $1/N$.  Choose their coordinate midpoints $y_j$ and put
$p_{k,N}(t)=\prod_{j=1}^k(t-N^2y_j)$.  Uniformly for $k/N$ in any compact
subinterval of $(0,\delta]$,
\begin{equation}\label{z3a:quantile-trial}
 |p_{k,N}(N^2y^2)|^2
 \le N^{4k+O(1)}\exp\{4NU_{b(v)}(y)\}\qquad(y\ge0).
\end{equation}
To check the uniform logarithmic error, the folded density bounds the
mass of an interval of length $h$ by $C\sqrt h(1+|\log h|)$.  Thus every
quantile cell has width at least $N^{-3}$ for large $N$.  Remove the cell
containing the evaluation point and at most two neighbors on each side.
On the remaining cells, monotonicity of the logarithmic kernel makes the
quadrature errors telescope, giving $O(\log N)$.  The removed cells have
total mass at most $5/N$; integration of logarithmic level sets, with
the preceding interval-mass bound, gives the same error.  Outside a fixed
compact interval the kernel has bounded total variation.  This proves
\eqref{z3a:quantile-trial}, including at polynomial zeros.

Monotone Riemann sums give
\[
 D_N(N^2y^2)\ge N^{2N}e^{NI_1(y)},\qquad
 C_{m_i}(N^2y^2)\le N^{2m_i}e^{NI_{s_i}(y)}.
\]
Combining these inequalities, the upper density bound for $\mu_3$, and
\eqref{z3a:equilibrium}, and then integrating, yields
\[
 \log\widehat h_{k,N}
 \le(4k+2m_1+2m_2-2N)\log N+N\ell(k/N)+O(\log N).
\]
The tail of the potential decreases at a uniform linear rate outside a
fixed compact interval, so all remaining powers in the density contribute
only $O(\log N)$.  Stirling's formula in \eqref{z3a:norm-product} cancels
all $N\log N$ terms.  It gives
\begin{equation}\label{z3a:analytic-rate}
 \limsup_{N\to\infty}N^{-2}\log|F_N(\xi)|\le-C_{\rm an},
 \qquad C_{\rm an}=\int_0^\delta A(v)\,\mathrm dv,
\end{equation}
where
\begin{align}
 A(v)={}&\sum_{i=1}^2\left[
  2(s_i+v)\log(4(s_i+v))
  -2s_i\log(s_i+\sqrt{s_i^2+b(v)^2})\right]\notag\\
 &-4v\log b(v)+2\log\frac{1+\sqrt{1+b(v)^2}}2,
 \qquad A(0)=2(s_1+s_2)\log2.\label{z3a:A}
\end{align}
For the omitted initial degrees $k<\epsilon N$, the monic trial $t^k$
suffices.  Bound its numerator by a power of $x+O(N)$, use
$D_N(x^2)\ge(N!)^2$, and integrate the exponential density.  The resulting
gamma-factor quotient is bounded by $e^{CN}N^C$ uniformly in those degrees.
Their total contribution is $O(\epsilon N^2)+O(N\log N)$; now let
$\epsilon\downarrow0$.  This completes the proof of
\eqref{z3a:analytic-rate}.

The monotonicity used for the finite certificate is also exact:
\[
 A'(v)=2\log\frac{16(s_1+v)(s_2+v)}{b(v)^2}>0.
\]
Indeed $v(b)\ge b/2$, so the logarithm's argument is at least $4$.
Left Riemann sums therefore give rigorous lower bounds for $C_{\rm an}$.

\subsection{A polynomial norm comparison, proved directly}

Set $z=t/N^2$ and normalize the physical quadratic form by the common
positive power of $N$:
\[
 \mathcal G_N(q)=N^{2N-2m_1-2m_2}
 \int_0^\infty |q(t/N^2)|^2
        \frac{C_{m_1}(t)C_{m_2}(t)}{D_N(t)}\,\mathrm d\mu_3(t).
\]
For a fixed $v>0$ and $\deg q\le vN$, one has, locally uniformly for
$z\notin[0,b(v)^2]$,
\begin{equation}\label{z3a:pointwise-norm}
 |q(z)|^2\le \mathcal G_N(q)
  \exp\{N(4U_{b(v)}(\sqrt z)-\ell(v))+o(N)\}.
\end{equation}
The error is independent of $q$; the square-root notation is unambiguous
because $\eta_b$ is symmetric.

We include the norm comparison needed to prove this statement.  On any
fixed compact interval in the positive $z$ axis, monotone product estimates
and the lower density bound imply that the density of $\mathcal G_N$ is
at least
\[
 cN\,\frac{z^{5/2}}{(z+s_1^2)(z+s_2^2)(z+1)}
          e^{-NV(\sqrt z)}\qquad(z>0).
\]
The exact constant is immaterial; the inequality follows by applying
the lower left-sum bound to each numerator product and the upper
right-sum bound to the denominator.  In particular the only loss near
zero is a fixed power of $z$.
Let $M=\max_{0\le z\le b^2}|q(z)|e^{-NV(\sqrt z)/2}$ and let $z_0$
be a maximizing point.  On a one-sided or two-sided interval of length
$N^{-1/2}$ about $z_0$, the $1/2$-H\"older continuity of $V(\sqrt z)$
makes $\sup|q|/|q(z_0)|\le e^{O(N^{3/4})}$.  The ordinary Markov
inequality on that interval then provides a subinterval of length
$e^{-O(N^{3/4})}N^{-O(1)}$, lying at least that distance from zero,
on which $|q|\ge|q(z_0)|/2$.  Integrating the preceding lower density
bound over that interval proves
\[
 M^2\le e^{o(N)}\mathcal G_N(q).
\]
This argument is uniform over all polynomials of degree at most $vN$.

To finish, push $\eta_{b(v)}$ to a measure $\sigma_v$ on $[0,b(v)^2]$.
Its mass is $v$ and its potential is $U_{\sigma_v}(z)=2U_{b(v)}(\sqrt z)$.
The function $\log|q(z)|-NU_{\sigma_v}(z)$ is subharmonic off that interval.
It has no positive logarithmic growth at infinity.  The contact identity
\eqref{z3a:equilibrium} and the maximum principle on the compactified
complement therefore give
\[
 \log|q(z)|-NU_{\sigma_v}(z)\le\log M-N\ell(v)/2.
\]
Combining the last two inequalities proves \eqref{z3a:pointwise-norm}.
Thus no asymptotic theorem about orthogonal polynomials is being assumed.

\subsection{Contour majorant and its angular deficit}

Fix $u>1$, put $a=u^2$, and use the circle $|z|=a$.  The $X$ coefficient
of $\mathcal J_X(R(t))$ is the sum of the residues of $-tR(t)$ at its
finite poles.  Consequently the coefficient bilinear form, in scaled
polynomial coordinates, is
\[
 -\frac{N^4}{2\pi i}\int_{|z|=a}
 z\frac{C_{m_1}(N^2z)C_{m_2}(N^2z)}{D_N(N^2z)}p(z)q(z)\,\mathrm dz.
\]
All poles are inside the circle.  Include the same positive normalization
as in $\mathcal G_N$, and take the absolute contour measure and the
Hermitian product $p\overline q$ to obtain a positive Gram form.
In a real basis orthonormal for the corresponding physical form let its
matrix be $\mathcal A_N$; the coefficient matrix in this basis is
orthogonally similar to $T_N$.  The common positive normalization cancels
in both relative forms.  Conjugation symmetry makes $\mathcal A_N$ real
symmetric.  Andreief's identity and
the triangle inequality applied to every principal index set give
\[
 |\det(T_N)_I|\le\det(\mathcal A_N)_I.
\]
Expanding by principal minors proves, for every complex $h$,
\begin{equation}\label{z3a:majorant}
 |\det(I+hT_N)|\le\det(I+|h|\mathcal A_N).
\end{equation}
Only this determinant inequality is used; an ordering of the absolute
eigenvalues of $T_N$ is unnecessary.

The exponential rate of the absolute contour density is
\[
 L(z)=\int_0^{s_1}\log|z+x^2|\,\mathrm dx
     +\int_0^{s_2}\log|z+x^2|\,\mathrm dx
     -\int_0^1\log|z+x^2|\,\mathrm dx.
\]
Its Riemann-sum error on this circle is $O(\log N)$, including all
polynomial prefactors.  Define
\[
 \Phi_v(z)=L(z)+4U_{b(v)}(\sqrt z)-\ell(v),\qquad
 \Gamma_u(v)=\Phi_v(-u^2).
\]
For $z=u^2e^{i\theta}$ put $y=1+\cos\theta\in[0,2]$.  The rational part
has deficit at least
\begin{equation}\label{z3a:rational-deficit}
 L(-u^2)-L(z)\ge\frac12\int_{s_1+s_2}^1
 \log\left(1+\frac{2u^2x^2}{(u^2-x^2)^2}y\right)\,\mathrm dx.
\end{equation}
Indeed the logarithmic integrand is increasing in $x$, and
$\int_0^{s_1}+\int_0^{s_2}$ is bounded above by $\int_0^{s_1+s_2}$.

The potential part has the further deficit
\[
 4U_{b(v)}(iu)-4U_{b(v)}(\sqrt z)
 =-\int\log\left(1-\frac{2u^2x}{(u^2+x)^2}y\right)
                 \,\mathrm d\sigma_v(x)
 \ge C_vy,
\]
where the arcsine mass derivative gives
\begin{equation}\label{z3a:linear}
 C_v=u\int_0^v\frac{b(w)^2}{(u^2+b(w)^2)^{3/2}}\,\mathrm dw.
\end{equation}
Both deficits are nonnegative.  Thus $\Phi_v$ is maximal at $-u^2$.
The same mass derivative gives the useful exact formula
\begin{equation}\label{z3a:Gamma}
 \Gamma_u'(v)=4\operatorname{arsinh}(u/b(v)),\qquad
 \Gamma_u(0)=I^-_{s_1}+I^-_{s_2}-I^-_1-\ell(0),
\end{equation}
where
\[
 I^-_s=s\log(u^2-s^2)+u\log\frac{u+s}{u-s}-2s.
\]

Choose any positive lower bound $c_\theta\le C_{\delta-T_{\rm cap}}$.
Partition $[s_1+s_2,1]$ into intervals of widths $2K_j$ with midpoints
$c_j$, and set
\[
 e_j=\frac{(u^2-c_j^2)^2}{2u^2c_j^2},\qquad
 D(y)=c_\theta y+\sum_jK_j\log(1+y/e_j).
\]
Then, simultaneously for $0\le t\le T_{\rm cap}$,
\begin{equation}\label{z3a:angular-majorant}
 \Phi_{\delta-t}(z)\le\Gamma_u(\delta-t)-D(y).
\end{equation}
For the midpoint assertion, after scaling $x$ by $u$, the function
$\log(1+2yx^2/(1-x^2)^2)$ is convex on $[0,1)$ for $0\le y\le2$.
After multiplication by positive factors its second derivative has
numerator
\[
 (1-y/2)(1-x^2)^2(3x^4+8x^2+1)
 +(y/2)(3x^8+6x^6+4x^4+2x^2+1)\ge0.
\]
Jensen's inequality on each interval proves the claimed lower sum.

\subsection{Distributed zeros and an eigenvalue envelope}

For $0\le h\le1$ define
$Q_h(z)=z^2+2a(1-h)z+a^2$.  On the physical axis $z\ge0$ it is at least
$a^2$, while on $|z|=a$ it satisfies
$|Q_h(z)|^2=a^4(2|y-h|)^2$.
For $k=2\ell$, restrict to polynomials
\[
 p(z)=\prod_{j=1}^{\ell}Q_{h_j}(z)q(z),\qquad \deg q<r-k.
\]
This is a real subspace of codimension $k$.  The lower physical factor
$a^{4\ell}$ cancels the same contour factor.  Equations
\eqref{z3a:pointwise-norm} and \eqref{z3a:angular-majorant} therefore
bound its contour-to-physical squared-norm ratio by the rate
\[
 \Gamma_u(\delta-t)+t\log2+
 2U_\rho(y)-D(y),\qquad k/N\longrightarrow t,
\]
if the empirical angular zeros converge to a measure $\rho$ of mass $t/2$.

The requisite measure is explicit.  Determine $B=B(t)$ by
\begin{equation}\label{z3a:angular-mass}
 t=\frac{c_\theta B}{2}+
       \sum_jK_j\left(1-\frac1{\sqrt{1+B/e_j}}\right),
\end{equation}
and suppose $B(t)\le1$.  Put $S_j=\sqrt{1+B/e_j}$.  The density
\[
 \frac{\mathrm d\rho}{\mathrm dy}
 =\frac1{2\pi}\sqrt{\frac{B-y}{y}}
     \left(c_\theta+\sum_j\frac{K_j}{S_j(y+e_j)}\right),
 \qquad0<y<B,
\]
has mass $t/2$, and $2U_\rho-D$ is constant on $[0,B]$ and no larger
on $[0,\infty)$.  Its Cauchy transform is
\[
 \frac12\left[c_\theta+\sum_j\frac{K_j}{z+e_j}
 -\sqrt{\frac{z-B}{z}}
     \left(c_\theta+\sum_j\frac{K_j}{S_j(z+e_j)}\right)\right].
\]
Cancellation at $-e_j$, its jump, and its coefficient at infinity prove
the transform and mass assertions.  The derivative of $2U_\rho-D$
outside $[0,B]$ is the negative square-root term in that transform,
which proves the global potential inequality.  Its contact constant,
after addition of $t\log2$, is
\begin{equation}\label{z3a:F}
 F(t)=t\log(B(t)/2)-\frac{c_\theta B(t)}2
       -2\sum_jK_j\log\frac{1+\sqrt{1+B(t)/e_j}}2.
\end{equation}
In particular $F(0)=0$ and $F'(t)=\log(B(t)/2)$.

Actual zero sets follow by equal-mass quantile discretization.  The
density here is bounded by $Cy^{-1/2}$, so the same logarithmic
quadrature argument used in \eqref{z3a:quantile-trial} gives a uniform
upper potential error $O(\log N/N)$ on $[0,2]$.  Choose
$k=2\lceil tN/2\rceil$; its rounding changes the rate by $o(1)$.  Only
this upper estimate is required.  Positive-matrix min--max now gives,
for each fixed $t\in[0,T_{\rm cap}]$, the $(k+1)$st eigenvalue of
$\mathcal A_N$ in decreasing order the exponential upper rate
\begin{equation}\label{z3a:envelope}
 g(t)=\Gamma_u(\delta-t)+F(t).
\end{equation}
This function is continuous and strictly decreasing while $B(t)\le1$,
because \eqref{z3a:Gamma} is positive and $F'(t)<0$.


\subsection{A closed expression used by the interval checker}

For reproducing $\Gamma_u$ by elementary directed evaluations instead of
an improper quadrature, put $U=\sqrt{u^2+b^2}$,
$A_s=\sqrt{s^2+b^2}$, and
\[
 J_s(b)=\log\frac{U+A_s}{u+s}
 -\frac{s}{u}\log\frac{uA_s+sU}{s(U+u)}.
\]
Then the exact formula used in the replay is
\[
 \Gamma_u(v(b))=\Gamma_u(0)
 +4v(b)\log\frac{u+U}{b}
 +2u\left[\log\frac{b+U}{u}-J_1(b)+J_{s_1}(b)+J_{s_2}(b)\right].
\]
Differentiation and the limit at $b=0$ verify it from
\eqref{z3a:Gamma}.  The only transcendental operations in this analytic
certificate are square roots and logarithms, enclosed outwards; the
arithmetic certificate and its interval-cell integrations are exact
rational calculations.

% Standalone proof core, no preamble. Requires amsmath, amssymb, amsthm.
% The host document should define lemma, proposition, theorem, corollary.
% All labels are prefixed zthree:.

\section{The cubic functional and factorial windows}
\label{zthree:setup}

Put $\xi=\zeta(3)$ and
\[
 C_m(t)=\prod_{a=0}^{m-1}(t+a^2),\qquad
 D_N(t)=\prod_{i=1}^N(t+i^2),\qquad
 b_j(t)=\frac{(-1)^jC_j(t)}{(2j)!},\qquad
 P_j(x)=b_j(-x^2).
\]
The rational functional used throughout this appendix is
\begin{align}
 \mathcal J_{3,X}(t^k)&=\frac{(-1)^k(2k+3)}2B_{2k+2},\label{zthree:moments}\\
 \mathcal J_{3,X}\left(\frac1{t+i^2}\right)
 &=i^2X-i^2H_{i-1}^{(3)}-\frac1{2i}-\frac12.\label{zthree:poles}
\end{align}
Polynomial division and partial fractions define it unambiguously on
rational functions with simple poles among the negative integer squares.
Here $B_1=-1/2$ and $H_n^{(3)}=\sum_{j=1}^n j^{-3}$.
At $X=\xi$ it is integration against the positive measure
\begin{equation}\label{zthree:measure}
 d\mu_3(x^2)=4\pi^2x^3\sum_{k\ge1}k^2e^{-2\pi kx}\,dx,
 \qquad x>0.
\end{equation}
The moment formula follows by termwise integration and Euler's even-zeta
formula. The corresponding Stieltjes identity is
\[
 \int\frac{d\mu_3(t)}{t+y^2}
   =y^2\zeta(3,y)-\frac1{2y}-\frac12\qquad(y>0).
\]
It follows, for example, by differentiating Binet's formula twice and
integrating by parts. In particular, \eqref{zthree:poles} at $X=\xi$
has the asserted positive-measure interpretation.

For $1\le m_1\le m_2\le N$, $r\ge1$, and $R_i=m_i+r$, define
\begin{equation}\label{zthree:matrix}
 A_{jk}(X)=(N!)^2\mathcal J_{3,X}
 \left(\frac{b_{m_1+j}b_{m_2+k}}{D_N}\right),\qquad 0\le j,k<r.
\end{equation}
Constant triangular row and column changes identify its determinant,
up to a nonzero rational factor, with that of the symmetric Hankel pencil
\begin{equation}\label{zthree:hankel}
 H_N(X)=\left[\mathcal J_{3,X}
 \left(\frac{C_{m_1}C_{m_2}t^{j+k}}{D_N}\right)\right]_{0\le j,k<r}.
\end{equation}
Its value $G_N=H_N(\xi)$ is positive definite. All valuations below are
coefficientwise Gauss valuations on rational polynomials in $X$.

\section{Local precision and the second arithmetic layer}
\label{zthree:local}

For the local statements take
\begin{equation}\label{zthree:good}
 p>5,\qquad p\le N,\qquad
 p^2>4\max(R_1,R_2)+20,\qquad p^2>2N+20.
\end{equation}
Write $N=qp+\sigma$, $0\le\sigma<p$, and put
\[
 \epsilon=\left\lfloor\frac{2\sigma}{p}\right\rfloor,
 \quad n=\frac{p-1}{2},\quad h=2q+1+\epsilon.
\]
Represent nonzero reflected classes by $1\le c\le n$. Their signed
physical levels are $K_c=\{a\in\mathbb Z:1\le|c+pa|\le N\}$.
There are $h$ levels on the maximal classes $M$ and $h-1$ on their
complement $T$, where
\begin{equation}\label{zthree:classes}
\begin{array}{c|c|c|c}
\epsilon&M&T&w=|M|\\\hline
0&\{1,\ldots,\sigma\}&\{\sigma+1,\ldots,n\}&\sigma\\
1&\{p-\sigma,\ldots,n\}&\{1,\ldots,p-\sigma-1\}&\sigma-n.
\end{array}
\end{equation}
The nonmaximal level sets are $\{-q,\ldots,q-1\}$ and
$\{-q,\ldots,q\}$, respectively.

\begin{lemma}[Precision after one derivative]\label{zthree:precision}
If $g$ is integer-valued at $p$, $\deg g=L<p^2$, and
$E=\lfloor L/p\rfloor$, then
\[
 v_p([Y^j]g(c+pY))\ge\max(0,j-E).
\]
If $E\le p-3$, reduction modulo $p$ or $p^2$ may precede one derivative
and Bernoulli evaluation. The respective retained degrees are $E$ and
$E+1$.
\end{lemma}
\begin{proof}
Expand in the integral Newton basis $\binom xa$. Among the numerator
factors $c+pY-i$, either $\lfloor a/p\rfloor$ or one more has constant
term divisible by $p$. Extracting these factors proves the coefficient
bound. For omitted terms of degree $j<p$, $B_{j-1}$ is integral.
For $j\ge p$, their coefficients have valuation at least three when
$E\le p-3$; the possible one-power Bernoulli denominator still leaves
a multiple of $p^2$. The modulo-$p$ argument is identical with one
less retained order. Constant terms disappear under differentiation.
\end{proof}

For window polynomials $f,g$ in the $x$ variable, let $I$ be the
factorial-basis truncation of $fg$ through index $N$, and set
\begin{equation}\label{zthree:quotient}
 Q_{fg}(x)=(N!)^2\frac{f(x)g(x)-I(x)}{D_N(-x^2)}.
\end{equation}
The numerator vanishes at $0,\pm1,\ldots,\pm N$. The relevant
Bernoulli input satisfies
\[
 \deg(x^3Q_{fg})\le2R_1+2R_2-2N-1
                  \le4\max(R_1,R_2)-2N-1.
\]
Since $N\ge p$, \eqref{zthree:good} implies $E\le p-3$ for this
input. Thus Lemma~\ref{zthree:precision} gives the precision needed
below. This is an explicit sufficient replacement for the
three-derivative precision argument. Alternatively, the first possible
pole $B_{p-1}$ is canceled by the first-derivative factor $p$; no
unproved precision gain is being assumed.

Define
\[
 Z_{2a}(Y)=\prod_{v=-a}^{a-1}(Y-v),\qquad
 Z_{2a+1}(Y)=\prod_{v=-a}^{a}(Y-v),\qquad Z_0=1.
\]
For $K=\{-q-\eta,\ldots,q\}$, $\eta\in\{0,1\}$, write
$L(Y)=\prod_{a\in K}(Y-a)$ and let $I_P$ be the remainder of $P$
modulo $L$. Set
\[
 h_3(a)=\begin{cases}\sum_{j=1}^a j^{-3},&a\ge0,\\
                         \sum_{j=1}^{-a-1}j^{-3},&a\le-1,
          \end{cases}
\]
and introduce the rational functional
\begin{equation}\label{zthree:universal}
 \mathcal T_{q,\eta}^{(3)}(P)=
 \sum_{a\in K}\frac{h_3(a)P(a)}{L'(a)}
 +\frac12\mathcal B\left[\left(\frac{P-I_P}{L}\right)'\right],
 \qquad \mathcal B(Y^j)=B_j.
\end{equation}
In the cases used here its required denominators are units at good
primes; alternatively, for each fixed limiting prime band one may
exclude a further finite set of primes.
For $S(Y)=-1-Y$ one has
\begin{equation}\label{zthree:parity}
 \mathcal T_{q,1}^{(3)}(P\circ S)
       =-\mathcal T_{q,1}^{(3)}(P).
\end{equation}
Indeed $K,L,h_3$ are invariant, $L'$ changes sign, and the quotient
derivative changes sign while $\mathcal B(F(-1-Y))=\mathcal B(F(Y))$.

Put
\begin{equation}\label{zthree:features}
 D_i=\left\lfloor\frac{2R_i-2}{p}\right\rfloor,\quad
 J_i=\lceil D_ip/2\rceil,\quad H_i=R_i-J_i,\quad
 k_i(c)=\#\{a:|c+pa|<m_i\}.
\end{equation}
When using complete features assume $m_i\le J_i$. Define
\begin{equation}\label{zthree:AK}
 L_i=\sum_{c\in M}(D_i-k_i(c))_+,\qquad
 A_i=L_i+\min(H_i,w),\qquad K_i=(H_i-w)_+,
 \quad x_+=\max(0,x).
\end{equation}

\begin{lemma}[Restriction dimensions]\label{zthree:dimensions}
The maximal restriction map on the $i$th window has rank $A_i$.
Its kernel has $K_i$ remaining top-coefficient directions. Its low
features on class $c$ are
$Z_{k_i(c)}\mathbb F_p[Y]_{\le D_i-1-k_i(c)}$.
\end{lemma}
\begin{proof}
The reduction of the span of $P_0,\ldots,P_{J_i-1}$ injects into the
direct sum of degree-$\le D_i-1$ polynomials on the nonzero classes
and even polynomials of that degree on the zero class. A zero image
means all integer values are divisible by $p$ and hence, by the
triangular factorial basis, all coefficients are divisible by $p$.
The target dimension is $D_i(p-1)/2+\lceil D_i/2\rceil=J_i$, so
the map is onto. The independent conditions at $0,\ldots,m_i-1$
give precisely the displayed roots, independently on each class.

Lucas reduction of the central-binomial formula for $2P_j$ shows
that the leading local coefficients of the top $H_i$ vectors are
polynomials of successive degrees $0,\ldots,H_i-1$ in $c^2$, with
a common extra factor $c$ if $D_i$ is odd. Their leading coefficients
are units. More explicitly the complete factors in the even and
odd bands are $P_a(c)\binom{Y+b}{2b}$ and
$cQ_a(c^2)\binom{Y+b}{2b+1}$, where
$Q_a(Z)=\prod_{v=1}^a(Z-(v-1/2)^2)/(2a+1)!$.
Evaluation at the $w$ distinct nonzero squared nodes has rank
$\min(H_i,w)$. Once these top coefficients vanish on $M$, low
features cancel the remaining restrictions independently. This
proves both assertions.
\end{proof}

For $u_1,u_2\in\{0,1\}$ put $n_i(c;u_i)=(D_i+u_i-k_i(c))_+$.
Use the cap $b_c(u_1,u_2)=\min(n_1,n_2)$, except that if
$\epsilon=0$ and both $k_i(c)$ are even use
\begin{equation}\label{zthree:paritycap}
 b_c(u_1,u_2)=
 \min(\lceil n_1/2\rceil,\lfloor n_2/2\rfloor)
 +\min(\lfloor n_1/2\rfloor,\lceil n_2/2\rceil).
\end{equation}
Finally set
\begin{equation}\label{zthree:B}
 B=\min\left\{r-A_1,r-A_2,
 \min_{u_1,u_2}\left(\sum_{c\in T}b_c(u_1,u_2)
                       +(1-u_1)K_1+(1-u_2)K_2\right)\right\}.
\end{equation}

\begin{theorem}[Cubic second layer]\label{zthree:second}
Under the hypotheses above,
\begin{equation}\label{zthree:secondbound}
 v_p(\det A(X))\ge-(1+\epsilon)r-A_1-A_2-B.
\end{equation}
Independently of the complete-feature hypothesis, the first layer gives
\begin{equation}\label{zthree:firstbound}
 v_p(\det A(X))\ge-(2+\epsilon)r
       -\min\{r,(\min(D_1,D_2)+1)w\}.
\end{equation}
\end{theorem}
\begin{proof}
Factor locally
$D_N(-(c+pY)^2)=p^{h_c}L_c(Y)U_c(Y)$, where $U_c$ is a constant
unit modulo $p$. At the zero class, extracting $p^2$ from each
physical multiple shows that $Q_{fg}(pY)$ is integral. Generally
$p^{1+\epsilon}Q_{fg}$ is integer-valued.

The Bernoulli representation and distribution are
\[
 \mathcal J_{3,X}(q)=\tfrac12\mathcal B\partial_x[x^3q(-x^2)],
 \qquad
 \mathcal B(F)=p^{-1}\sum_{c=0}^{p-1}\mathcal B_Y F(c+pY).
\]
Partial fractions and these identities show that
$p^{3+\epsilon}A\bmod p$ factors through maximal restrictions,
with nonzero multipliers
$\omega_c=2c^3(N!/p^q)^2/\overline U_c$ and local functional
\eqref{zthree:universal}. This proves \eqref{zthree:firstbound} by
factoring from restriction-kernel rows. The coefficient of $X$ in
$A$ has valuation at least $-\epsilon$, so the layers in question
are constant in $X$.

On both maximal kernels $fg$ is locally divisible by $p^2$ on $M$.
The reduction of $I$ has degree at most $h-1$ and $h$ physical zeros,
so vanishes. The reduction of $I/p$ has degree at most $h$ and is
a multiple of $L_c$. Thus $p^\epsilon Q_{fg}$ is constant modulo
$p$ on $M$. On $T$, $I$ differs from the interpolation remainder
by a constant multiple of $L_c$. Exactly,
\[
 p^{2+\epsilon}\operatorname{Quot}(f,g)
 =\tfrac12\sum_c\mathcal B_Y\partial_Y
       [(c+pY)^3p^\epsilon Q_{fg}(c+pY)].
\]
Lemma~\ref{zthree:precision} permits reduction first. Maximal and
interpolation constants disappear after differentiation. The zero
class disappears because of $(pY)^3$, even allowing a Bernoulli
denominator. Maximal proper fractions gain $p^2$ and disappear;
nonmaximal fractions supply the residue term. Physical multiples
of $p$ and the retained constant quotient from $P_N$ also disappear.
The surviving second form is the pullback of
$\bigoplus_{c\in T}\omega_c\mathcal T_*^{(3)}$, with
$\mathcal T_*^{(3)}=\mathcal T_{q-1,1}^{(3)}$ if $\epsilon=0$ and
$\mathcal T_{q,0}^{(3)}$ otherwise.

The feature dimensions bound its rank by \eqref{zthree:B};
\eqref{zthree:paritycap} follows from the two parity eigenspaces and
\eqref{zthree:parity}. In bases adapted to the two maximal kernels,
all kernel rows and columns of $M_0=p^{3+\epsilon}A$ are divisible
by $p$, and their common block after division by $p$ has rank at
most $B$. Factoring gives
$v_p\det M_0\ge2r-A_1-A_2-B$, proving \eqref{zthree:secondbound}.
If the indicated extra row count is negative, a single-side
factorization gives the same weaker bound.
\end{proof}

\section{The upper prime band and small-prime clearance}
\label{zthree:monomial}

The exact factorial change has opposite valuation
\begin{equation}\label{zthree:factorial}
 f_p=\sum_{i=1}^2\sum_{j=m_i}^{R_i-1}\left\lfloor\frac{2j}{p}\right\rfloor
            -2r\left\lfloor\frac Np\right\rfloor
\end{equation}
at good primes. For the monomial matrix \eqref{zthree:hankel}, division
gives a quotient matrix plus rank-one residue atoms. Its fractional
quotient part is supported in a terminal block of size
\begin{equation}\label{zthree:terminal}
 t_p=\min\{r,(R_1+R_2-N-p+1)_+\}.
\end{equation}
Indeed the first possible moment pole, $2k+2=p-1$, is canceled by
$2k+3=p$ in \eqref{zthree:moments}; the first uncanceled possible
index is $k=p-2$. Each needed moment costs at most one power of $p$.
A fractional entry consequently requires
$j+k\ge N+p-2-m_1-m_2$, forcing both indices into the stated block.
Every quotient minor of order $a$ costs at most $\min(a,t_p)$.

\begin{proposition}[Cubic monomial bound]\label{zthree:monobound}
Assume $N/3<p\le N$, $p>2\max(m_1,m_2)$, the square-size
hypotheses in \eqref{zthree:good}, and
$r\ge H_p=N-p-(q-1)$. For $p>N/2$ put
$a_m(p)=\min(m-1,\sigma)+(m-p+\sigma)_+$; for $p<N/2$ put
$a_m(p)=0$. Then
\[
 v_p\det A\ge-f_p-5H_p+a_{m_1}(p)+a_{m_2}(p)
                         -\min(t_p,r-H_p).
\]
For $p>N$ satisfying the square-size hypotheses, all residue
contributions are integral, and the valid bound is
$v_p\det A\ge-f_p-t_p$.
\end{proposition}
\begin{proof}
In a nonzero reflected class with two low nodes and $h$ high nodes,
a mixed-minor term selecting $\ell$ atoms, $b$ of them high, costs
at most $3b+\ell(h+2-\ell)$. The harmonic loss is three; each
derivative denominator costs $h+1$; the two evaluation minors restore
$\ell(\ell-1)$ powers through their Vandermonde factors. For
$1\le h\le4$ the maximum is $5h$. Omitting $d$ atoms below $h$
loses at least $d(5-h)+d^2\ge d$ from this maximum.

When $p>N/2$, one has $h\le2$. If $a\in\{0,1,2\}$ common
window factors affect a class, surviving atoms gain $a$ powers and
the low root node vanishes. The maximum becomes $(5-a)h$; the
deficit for $d$ omitted atoms is $d(5-a-h)+d^2\ge d$. Each window
therefore saves its own affected high-node count $a_m(p)$.
For $N/3<p<N/2$ we deliberately discard all such savings: with
four high nodes and two common factors the one-unit deficit need
not hold. The unaffected bound remains valid there.

Unaffected all-low pairs are integral by the reflection identity
$k_{p-i}^{(3)}\equiv k_i^{(3)}\pmod p$, which follows from
$\sum_{a=1}^{p-1}a^{-3}=0$. Affected low atoms are integral as
well. In the zero class both common factors contribute valuation
two, the cubic pole scalar costs at most one, and the derivative
costs $2(q-1)$; its atoms have valuation at least $5-2q\ge1$.
Absorb these integral pieces into the quotient's integral part.

If $\ell$ atoms remain in a mixed-minor term, their cost is at most
$5H_p-a_{m_1}-a_{m_2}-(H_p-\ell)_+$. The quotient minor costs
at most $\min(t_p,r-\ell)$. Since $r\ge H_p$,
\[
 \min(t_p,r-\ell)-(H_p-\ell)_+\le\min(t_p,r-H_p).
\]
This proves the first bound after restoring \eqref{zthree:factorial}.
For $p>N$, every residue class is a low singleton or reflected pair,
so the same cancellation gives the second assertion.
\end{proof}

\begin{lemma}[Small primes]\label{zthree:small}
For every prime, with $R=\max(R_1,R_2)$ and
$L_p=\lfloor\log_p(4R+4)\rfloor$, each entry satisfies
\[
 v_p(A_{jk})\ge-9L_p-v_p(96).
\]
If all parameters are $O(N)$, clearing primes $p\le N^{9/10}$
in the determinant costs $o(N^2)$ logarithmically.
\end{lemma}
\begin{proof}
Expand products in the even integer-valued factorial basis; at $2$
this loses at most two fixed powers. For basis index $l\le N$, use
\[
 D_N'(-i^2)=\frac{(-1)^{i-1}(N-i)!(N+i)!}{2i^2},\qquad
 \frac{(N!)^2}{(N-i)!(N+i)!}
   =\frac{\binom{2N}{N-i}}{\binom{2N}{N}}.
\]
The ratio costs at most $L_p$, and the cubic harmonic term at most
$3L_p$ plus fixed powers of two. The constant quotient at $l=N$
has moment $1/4$ and is covered by the stated bound.
For $l\ge N+1$, put $k=l-N-1$. The quotient's two products of
$k$ consecutive integers are divisible by $(k!)^2$, and
\[
 \frac{(N!)^2(k!)^2}{(2l)!}
   =\frac1{l^2\binom{2l}{l}\binom{l-1}{N}^2}
\]
costs at most $5L_p$. A derivative and Bernoulli evaluation each
cost at most $L_p$ in the Newton basis. These estimates are weaker
than $9L_p+v_p(96)$ and prove the assertion. The total log cost is
$O(rN^{9/10}\log N)=o(N^2)$.
\end{proof}

\section{The limiting arithmetic profile and its certificate}
\label{zthree:profile}

Take $m_i=s_iN$, $r=\delta N$, with
$0<s_1\le s_2<1/2$ and $\delta>1/2$ fixed, and
$\rho_i=s_i+\delta$. For $0<\beta<1$ put
\[
 q=\lfloor1/\beta\rfloor,\quad
 e=\lfloor2/\beta\rfloor-2q,\quad D_i=\lfloor2\rho_i/\beta\rfloor,
 \quad H_i=\rho_i-D_i\beta/2,\quad w=1-(q+e/2)\beta.
\]
Set $M_\beta=(0,w)$ if $e=0$ and
$M_\beta=(\beta/2-w,\beta/2)$ if $e=1$, and let
$T_\beta=(0,\beta/2)\setminus M_\beta$.
For $a_i=\lfloor s_i/\beta\rfloor$, $v_i=s_i-a_i\beta$, and
$z_i=\min(v_i,\beta-v_i)$, the normalized root count is
\[
 k_i(z)=\begin{cases}
 2a_i+1,&0<z<z_i,\\
 2a_i,&z_i<z<\beta/2,\ v_i<\beta/2,\\
 2a_i+2,&z_i<z<\beta/2,\ v_i>\beta/2.
 \end{cases}
\]
Values on the finite boundaries are immaterial. Define
\[
 \mathsf A_i=\int_{M_\beta}(D_i-k_i(z))_+\,dz+\min(H_i,w),
 \qquad \mathsf K_i=(H_i-w)_+.
\]
Define $b_z(u_1,u_2)$ by the dimension and parity rules
\eqref{zthree:paritycap}, with $n_i=(D_i+u_i-k_i(z))_+$, and put
\[
 \mathsf B=\min\left\{\delta-\mathsf A_1,\delta-\mathsf A_2,
 \min_{u_1,u_2}\left(\int_{T_\beta}b_z(u_1,u_2)\,dz
 +(1-u_1)\mathsf K_1+(1-u_2)\mathsf K_2\right)\right\}.
\]
The first and second local profiles are
\[
 E_1=(2+e)\delta+\min\{\delta,(\min(D_1,D_2)+1)w\},\qquad
 E_2=(1+e)\delta+\mathsf A_1+\mathsf A_2+\mathsf B.
\]
The complete-feature condition holds on $\beta<1$ because
$D_i\beta/2>\rho_i-\beta/2>s_i$.

For the monomial profile put
\[
 f(\beta)=\sum_{i=1}^2\left(D_i\rho_i-
           \frac{\beta D_i(D_i+1)}4\right)-2\delta q,
 \qquad t(\beta)=(\rho_1+\rho_2-1-\beta)_+,
\]
and set $a_i(\beta)=0$ below $1/2$ and
$a_i(\beta)=\min(s_i,1-\beta)+(1+s_i-2\beta)_+$ above $1/2$.
On the domain $\beta>\max(1/3,2s_2)$ and $\delta\ge1-\beta$, use
\[
 E_m=f(\beta)+5(1-\beta)-a_1(\beta)-a_2(\beta)
                       +\min\{t(\beta),\delta-1+\beta\}.
\]
For $\beta\ge1$ use $E_m=f(\beta)+t(\beta)$, with $q=0$.
Let $E$ be the pointwise minimum of available $E_1,E_2,E_m$ below
one and equal to $E_m$ above one. It vanishes beyond $2\rho_2$.

The prime number theorem, first on each fixed band away from zero,
converts these local bounds into a rational normalization whose rate is
at most $\int E(\beta)\,d\beta$. Negative local costs are allowed:
they record common factors of all coefficients. Lemma~\ref{zthree:small}
clears the small primes separately. Since $E_1\le4\delta$, for every
fixed $Q$ the following conservative bound is valid:
\begin{equation}\label{zthree:car}
 C_{\rm ar}\le\frac{4\delta}{Q}
                +\int_{1/Q}^{2\rho_2}E(\beta)\,d\beta.
\end{equation}
To justify the tail, clear $p\le N^{9/10}$ first, then use $4r$
for $N^{9/10}<p<N/Q$. All remaining primes are good for large $N$.
Fixed endpoints and integer rounding contribute $o(N^2)$.

For
\begin{equation}\label{zthree:parameters}
 (s_1,s_2,\delta)=\left(\frac2{25},\frac8{25},\frac{51}{50}\right),
 \qquad Q=500,
\end{equation}
exact rational cell integration gives
\begin{equation}\label{zthree:arcertificate}
 C_{\rm ar}\le
 \frac{1926282169226183949}{500000000000000000}
 =3.852564338452367898.
\end{equation}
The certificate enumerates every floor breakpoint and splits every
affine min/max crossing exactly. Its 6752 cells cover
$[1/500,67/25]$; each cell area is rounded upward to a multiple of
$10^{-18}$. The finite integral upper bound is
$1922202169226183949/500000000000000000$, and the tail is $51/6250$.
These rational operations, rather than floating quadrature, certify
\eqref{zthree:arcertificate}.

\section{Cubic Pearson identity and bounded root multiplicity}
\label{zthree:nullity}

Define $\mathcal K_X$ on polynomials by
$\mathcal K_X(f)=\mathcal B(f')/2$, and at simple integer poles by
\[
 \mathcal K_X\left(\frac1{x-i}\right)=
 \begin{cases}-X+H_i^{(3)},&i\ge0,\\
 -X+H_{-i-1}^{(3)},&i\le-1.
 \end{cases}
\]
The Bernoulli shift formula and direct evaluation at each pole give
\begin{equation}\label{zthree:K}
 \mathcal K_X(f(x+1)-f(x))=[x^2]f,\qquad
 \mathcal K_X(f)=-\tfrac12[x^2]f\quad(f\text{ even}),\qquad
 \mathcal J_{3,X}(h)=\mathcal K_X(x^3h(-x^2)).
\end{equation}
Put $W=x^3C_{m_1}(-x^2)C_{m_2}(-x^2)/D_N(-x^2)$ and
$\mathcal L_X(f)=\mathcal K_X(Wf)$. The function $W$ vanishes to
order seven at zero. With
\[
 A(x)=(x-1)^6(x-m_1)(x-m_2)(x+N),\quad
 B(x)=(x+1)^6(x+m_1)(x+m_2)(x-N),
\]
one has $B(x)=-A(-x)$ and $W(x+1)A(x+1)=W(x)B(x)$.
Equation~\eqref{zthree:K} therefore gives the exact Pearson identity
\begin{equation}\label{zthree:pearson}
 \mathcal L_X(B(x)f(x+1)-A(x)f(x))=0.
\end{equation}
The functional also vanishes on odd polynomials.

\begin{lemma}[Hankel kernel and polynomial adjoint]\label{zthree:kernel}
If a finite $r$ by $r$ Hankel matrix for a functional $L$ has nullity
$m>0$, its least-degree monic kernel polynomial $p$, of degree $j$,
satisfies $j\le r-m$, the leading $j$ by $j$ Hankel matrix is
nonsingular, and $L(pt^k)=0$ for $k\le r+m-2$.
If polynomial operators $S,S^*$ are adjoints for $L$ and each raises
degree by at most $d$, then orthogonality through degree $M$ and
$j+d-1\le M$ imply $p\mid Sp$.
\end{lemma}
\begin{proof}
Division by $p$ shows that a polynomial of degree below $j$
orthogonal to all smaller degrees would be in the full kernel,
contradicting minimality. This proves nonsingularity. Dividing any
kernel vector by $p$ then shows that its remainder vanishes.
Write $T_k=L(pt^k)$; initially $T_0,\ldots,T_{r-1}=0$.
If the first later nonzero index is $\nu$, the permissible multiples
of $p$ have degree at most $\nu-r$, so $m=\nu-r+1$.
If none of the relevant indices is nonzero, $m=r-j$.
Both cases prove the stated orthogonality range.
For the adjoint assertion, divide $Sp=pV+R$, $\deg R<j$,
$\deg V\le d$. Testing $R$ against $1,t,\ldots,t^{j-1}$ gives
zero by adjointness and the degree hypothesis. Nonsingularity forces
$R=0$.
\end{proof}

\begin{theorem}[Uniform cubic multiplicity bound]\label{zthree:multiplicity}
If $1\le m_1\le m_2\le N$ and $r>N-m_2+1$, then
$\operatorname{nullity}H_N(X)\le2$ for every complex $X$.
Every finite zero of its determinant has multiplicity at most two.
\end{theorem}
\begin{proof}
Define
\[
 \mathcal Df=B(x)(f(x+1)-f(x))-B(-x)(f(x-1)-f(x)),
 \quad \mathcal Qf=\operatorname{even}[B(x)f(x+1)],
 \quad \mathcal Rf=x\mathcal Df.
\]
Pearson makes $\mathcal D$ self-adjoint; on even polynomials the
adjoints of $\mathcal Q,\mathcal R$ are
$\operatorname{even}[A(x)g(x-1)]$ and $\mathcal D(xg)$.
Because $B$ is monic of odd degree nine, the degree-eight gain of
$\mathcal D$ cancels, leaving gain at most seven. All four even
operators have $x$-degree gain at most eight, hence $t$-degree gain
at most four under $t=-x^2$.

If the nullity were $m\ge3$, Lemma~\ref{zthree:kernel} would give
a monic $p$ of degree $j\le r-m$, orthogonal through $M=r+m-2$.
Since $j+3\le r-m+3\le M$, putting $P(x)=p(-x^2)$ gives
$P\mid\mathcal QP$ and $P\mid\mathcal RP$.
At a nonzero root $\alpha$ of $P$ the resulting sum and difference
imply $B(\alpha)P(\alpha+1)=0$. A noninteger root cannot reach
an integer root of $B$ or zero under repeated addition of one,
so would produce infinitely many roots. Thus all roots of $P$
are integers and all roots of $p$ are nonpositive squares.
Consequently $p(t)>0$ on $t>0$.

Now $E(t)=\prod_{i=m_2}^N(t+i^2)$ has degree below $r$, whereas
\[
 L_X(pE)=\mathcal J_{3,X}(p(t)tC_{m_1}(t))
       =\int p(t)tC_{m_1}(t)\,d\mu_3(t)>0.
\]
This contradicts the kernel equation, even for complex $X$.
Finally, congruence by $G_N^{-1/2}$ turns the pencil into
$I+(X-\xi)T_N$ with $T_N$ real symmetric. Its finite-root
multiplicities equal its nullities. Constant basis changes and
primitive normalization do not affect them.
\end{proof}

For \eqref{zthree:parameters} take $N$ through multiples of 50.
The slope matrix has exactly $N-m_2+1$ nonzero residue atoms;
their evaluation matrix has full column rank. Thus
\begin{equation}\label{zthree:degree}
 \deg\det H_N=N-m_2+1=\frac{17}{25}N+1,
 \qquad \lambda=\frac{17}{25}.
\end{equation}
The hypotheses of Theorem~\ref{zthree:multiplicity} hold.

\section{Determinant response and the quantitative transfer}
\label{zthree:transfer}

Let $P_N\in\mathbb Z[X]$ be the primitively normalized determinant,
$M_N=[X]H_N$, and $T_N=G_N^{-1/2}M_NG_N^{-1/2}$. Then
\[
 P_N(\xi+h)=P_N(\xi)\det(I+hT_N).
\]
The cubic slope is the residue-contour functional with numerator
factor $-t$, since its pole coefficient is $i^2=-t_i$.
In a real $G_N$-orthonormal basis, replace its contour measure by
its absolute value and conjugate the second polynomial to obtain
a positive matrix $\mathcal A_N$. Andreief's identity and the
triangle inequality give, for each principal index set $I$,
\[
 |\det(T_N)_I|\le\det(\mathcal A_N)_I,
 \qquad |\det(I+hT_N)|\le\det(I+|h|\mathcal A_N).
\]
This is determinant domination, not an asserted ordering of the
individual absolute eigenvalues of $T_N$.
The cubic physical density is positive, tends to $1/\pi$ at zero,
and has the same exponential tail as the quintic density. In fact
their exact ratio is
\[
 \frac{\rho_5(x)}{\rho_3(x)}
   =\frac{\pi^2x^2}{3}+\left(\frac{\pi x}{\sinh(\pi x)}\right)^2.
\]
Thus the weighted-polynomial arguments use the same leading
external field; the fixed polynomial factors are retained in the
subexponential estimates. The analytic certificate in the main text
supplies the required norm and response bounds for this cubic measure.

\begin{theorem}[Bounded-multiplicity transfer]\label{zthree:transferthm}
Suppose $\deg P_N\le\lambda N+O(1)$,
$\log|P_N(\xi)|\le-\varepsilon N^2+o(N^2)$, and rational root
multiplicities are uniformly bounded. If for some $a>0$ and every
fixed $C>0$,
\[
 \limsup_{N\to\infty}N^{-2}\log
 \sup_{|h|\le Ce^{-aN}}|\det(I+hT_N)|\le R(a)<\varepsilon,
\]
then $\mu(\xi)\le\lambda a/(\varepsilon-R(a))$.
\end{theorem}
\begin{proof}
For infinitely many rational approximants $x=p/q$ with
$|x-\xi|\le q^{-M}$, choose stages $N\sim c\log q$ with
$\lambda/(\varepsilon-R(a))<c<M/a$. A fixed stage modulus only
changes $N$ by $O(1)$. If $m=O(1)$ is the multiplicity of $x$ as
a root, with $m=0$ when it is not a root, integrality gives
$|P_N^{(m)}(x)/m!|\ge q^{-(\deg P_N-m)}$.
Cauchy's estimate on radius $e^{-aN}$ gives the upper logarithm
$-(\varepsilon-R(a))N^2+amN+o(N^2)$, contradicting the choice
of $c$. This excludes every $M>\lambda a/(\varepsilon-R(a))$.
\end{proof}

If a positive decreasing envelope $g(t)$ bounds the ordered
positive-contour eigenvalue exponents for $0\le t\le T<\delta$,
then, with $a=g(T)>0$ and $I(T)=\int_0^Tg(t)\,dt$,
\[
 R(a)\le I(T)-Ta,\qquad
 \mu(\xi)\le\frac{\lambda a}{\varepsilon-I(T)+Ta}.
\]
A finite partition and min--max bound the fixed-index blocks and
the remaining eigenvalues; refining the partition gives this
response estimate. A positive continuous pointwise envelope may
first be replaced by its running minimum. Its integral can only
decrease, so no unverified monotonicity assumption is necessary.

For $T=17/125$, the separate analytic and arithmetic certificates
give
\[
 \varepsilon-I(T)\ge S=0.004682016878632102,
 \qquad 0<a\le a_{\rm up}=6.545408942957.
\]
All displayed finite decimals are exact rational bounds. Since
$\lambda a/(S+Ta)$ is increasing in $a$ when $S>0$, the final
transfer is
\begin{equation}\label{zthree:finaltransfer}
 \mu(\zeta(3))\le
 \frac{\lambda a_{\rm up}}{S+Ta_{\rm up}}
 =\frac{2225439040605380000}{447428816560392051}
 =4.973839319768086\ldots<4.974.
\end{equation}
The positive endpoint and strict gap in this last step are inputs
from that analytic certificate; the present appendix supplies the
local arithmetic, multiplicity, and transfer justifications.

\section{Why a scalar retuning did not suffice}\label{sec:scalar}
A separate examination of the ordinary Pr\'evost scalar remainder Pad\'e family helps identify where the improvement comes from. Let the Pad\'e index be $m\sim rn$ and use the usual least-common-multiple denominator estimate. Define
\[
 L(r)=2\arsinh r+2r\arsinh(1/r),\qquad D(r)=3\max(1,r).
\]
In the range $L(r)>D(r)$, that standard certificate gives
\[
 U(r)=\frac{2L(r)}{L(r)-D(r)}.
\]
Since
\[
 L'(r)=2\arsinh(1/r)>0,\qquad L(r)=rL(1/r),
\]
$U$ decreases up to $r=1$ and satisfies $U(r)=U(1/r)$. Its minimum is therefore
\[
 U(1)=\frac{8\log(1+\sqrt2)}{4\log(1+\sqrt2)-3}
 =13.4178202333537\ldots.
\]
This is a limitation of the standard scalar certificate with its usual denominator bound. It does not rule out improved scalar divisibility, a different approximation family, or the growing determinants used in this report.

In the determinant construction, asymmetry changes the relevant degree rate from approximately one to $1-s_2$. The lower harmonic order also changes the local arithmetic exponent: the quintic second-layer term $(3+\epsilon)\delta$ becomes $(1+\epsilon)\delta$ after the cubic precision argument. The remaining rank terms, the larger cubic quotient-support cost, and the upper-prime contribution must all be included. A blanket subtraction of $2\delta$ from an old total arithmetic constant would be invalid. The exact piecewise profile is the reason the current result has a reliable certificate.

\section{Historical comparison and provenance}\label{sec:history}
The public benchmark located in the research search is Leibtag, Barkan, Kalisch, Weinbaum, Shalyt and Kaminer's October 2026 result $\mu(\zeta(3))<5.4197$~\cite{leibtag}. Their released certificates give a more precise comparison value about $5.419648258155345$; unpacking those digits would not itself be an improvement. Giraudin's September 2026 bound was $5.5138800$~\cite{giraudin}, following Rhin--Viola's $5.513891$~\cite{rv}. Litsey's recent L\"owner-determinant paper records $6.0382$ for $\zeta(3)$ and explicitly compares it with the stronger $5.4197$ benchmark~\cite{litsey}.
\begin{center}
\begin{tabular}{@{}lll@{}}
\toprule
Work & Upper bound & Status in this comparison\\
\midrule
Rhin--Viola (2001) & $5.513891$ & Earlier record\\
Giraudin (September 2026) & $5.5138800$ & Intermediate improvement\\
Leibtag et al. (October 2026) & $5.4197$ & Published preprint benchmark\\
This project, current report & $4.974$ & New research proof and certificates\\
\bottomrule
\end{tabular}
\end{center}
The claim of improvement is the numerical comparison supported by the proof in this report. A literature search cannot establish an absolute universal priority statement. The scalar functional is attributed to the Pr\'evost framework~\cite{prevost}; the current contribution is the quantitative asymmetric determinant argument and its specific certified exponent.

The local project sources consulted were the November 2026 versions of \path{zeta5_journal_publication.tex} and \path{zeta5_spectral_transfer.tex}. The first supplies the positive-measure and factorial-window framework, including a cubic specialization. The second develops the two-window arithmetic, bounded-root-multiplicity transfer, and distributed-zero contour method. The current report supplies the cubic adaptations, the asymmetric rational parameter choice, and the new numerical certificate. Both earlier sources are included as provenance in the reproducibility bundle, with their original contents preserved. The present bound does not assume an irrationality theorem for $\zeta(5)$ or an irrationality-measure bound for it.

\section{Reproducibility and limits of the certification}\label{sec:repro}
Unpack the accompanying archive and run
\begin{verbatim}
python verify.py
\end{verbatim}
The launcher first recomputes the rational arithmetic cells and their hash, then reruns the directed analytic checker and tests the exact final rational inequality. The finite certificate requires Python, NumPy, SciPy, and mpmath. The recorded environment used Python 3.12.14, NumPy 2.3.5, SciPy 1.17.0, and mpmath 1.3.0. SciPy's root solver only proposes brackets; no unchecked floating-point quadrature or root estimate contributes to a bound. The original exploratory determinant computations used SymPy, but SymPy is not needed to replay the theorem's certificate.

The main files in the archive are:
\begin{itemize}
\item \texttt{exact\_arithmetic.py}: exact rational profile construction and integration;
\item \texttt{certify\_analytic.py}: directed analytic and response inequalities;
\item \texttt{verify.py}: one-command replay and final exact comparison;
\item \texttt{certificates/}: the arithmetic cells, recorded numerical bounds, and environment;
\item \texttt{audits/}: separate arithmetic and multiplicity/transfer reviews;
\item \texttt{project\_sources/}: the earlier proof sources used for provenance;
\item the complete source and PDF of this report.
\end{itemize}
The arithmetic cell file has SHA256
\begin{center}\small\ttfamily
fcfa9c5d70e7f54d40c74296354a97b82b4a04b5be7db9afaa06ebe09878d049.
\end{center}
The finite check certifies the numerical inputs to the mathematical lemmas. Its successful execution does not replace the asymptotic proofs, the prime number theorem, or the functional identities. This work has received separate internal checks of the cubic arithmetic, asymmetric analytic formulas, interval-enclosure directions, and root-multiplicity transfer. It has not yet been externally refereed or formalized in Lean. The strongest asserted result is the ordinary theorem with its replayable certificates stated at the start of this report.

\begin{thebibliography}{9}
\bibitem{prevost}
M. Pr\'evost, \emph{Remainder Pad\'e approximants for the Hurwitz zeta function}, arXiv:1709.05389 (2017).
\url{https://arxiv.org/abs/1709.05389}.
The paper discusses the author's earlier remainder-Pad\'e irrationality proofs and their identification with Ap\'ery's sequences.
\bibitem{rv}
G. Rhin and C. Viola, \emph{The group structure for $\zeta(3)$}, Acta Arithmetica \textbf{97} (2001), 269--293.
\bibitem{giraudin}
D. Niedbala Giraudin, \emph{A new upper bound for the irrationality exponent of $\zeta(3)$}, arXiv:2609.26980 (22 September 2026).
\url{https://arxiv.org/abs/2609.26980}.
\bibitem{leibtag}
E. Leibtag, H. Barkan, R. Kalisch, S. Weinbaum, M. Shalyt and I. Kaminer, \emph{New Directions in Irrationality}, arXiv:2610.27799 (20 October 2026).
\url{https://arxiv.org/abs/2610.27799}.
Computational supplement:
\url{https://github.com/RamanujanMachine/new-directions-in-irrationality}.
\bibitem{litsey}
S. Litsey, \emph{Irrationality measures from L\"owner determinants}, hexagon:2610.00401v3 (manuscript dated 10 November 2026).
\url{https://hexagonmath.org/pdf/2610.00401v3}.
\end{thebibliography}
\end{document}
